\subsection{Examples of Gaussian promeasures}

\begin{enumerate}
\def\labelenumi{\arabic{enumi})}
\tightlist
\item
  Let E be a real Hilbert space. The mapping \(x' \mapsto \|x'\|^2\) is a positive quadratic form on \(E'\) . The corresponding Gaussian promeasure is called the canonical Gaussian promeasure on E. It can be shown that this promeasure is not a measure if E is infinite-dimensional.
\end{enumerate}

Let \(A\) be a continuous linear operator on \(\mathbf{E}\) . The mapping \(x' \mapsto \|{}^{t} A \cdot x'\|^2\) is a positive quadratic form on \(\mathbf{E}'\) . The corresponding promeasure \(\mu_A\) on \(\mathbf{E}\) is a measure if and only if \(A\) is a Hilbert--Schmidt operator (cf.~No.~11, Cor. 2 of Th. 3).

\begin{enumerate}
\def\labelenumi{\arabic{enumi})}
\setcounter{enumi}{1}
\tightlist
\item
  Kernels of positive type. Let T be a set and \(E = R^{T}\) the vector space of real functions on T, equipped with the topology of pointwise convergence. For every \(t \in T\) , one denotes by \(\varepsilon_{t}\) the linear form \(f \mapsto f(t)\) on E. The family \((\varepsilon_{t})_{t \in \mathbb{T}}\) is a basis of \(E'\) (TVS, II, §6, No.~6, Cor. 2 of Prop. 8).
\end{enumerate}

One calls (real) kernel of positive type on T every real-valued function K on \(T \times T\) satisfying the relations

\[
\mathbf {K} (t, t ^ {\prime}) = \mathbf {K} (t ^ {\prime}, t) \quad \text { for } t, t ^ {\prime} \text { in } \mathbf {T},\tag{20}
\]

\[
\sum_ {i, j = 1} ^ {p} c _ {i} c _ {j} \mathrm{K} (t _ {i}, t _ {j}) \geqslant 0\tag{21}
\]

for any positive integer \(p\) , elements \(t_1, \ldots, t_p\) of \(T\) , and real numbers \(c_1, \ldots, c_p\) . If this is so, the formula

\[
q \left(\sum_ {t \in \mathrm{T}} c _ {t} \varepsilon_ {t}\right) = \sum_ {t, t ^ {\prime} \in \mathrm{T}} c _ {t} c _ {t ^ {\prime}} \mathrm{K} (t, t ^ {\prime})\tag{22}
\]

defines a positive quadratic form on \(\mathbf{E}'\) . Conversely, if \(q\) is a positive quadratic form on \(\mathbf{E}'\) , then the formula

\[
\mathbf {K} (t, t ^ {\prime}) = \frac {1}{2} [ q (\varepsilon_ {t} + \varepsilon_ {t ^ {\prime}}) - q (\varepsilon_ {t}) - q (\varepsilon_ {t ^ {\prime}}) ]\tag{23}
\]

defines a kernel K of positive type on T. One thus obtains two mutually inverse bijections between the set of kernels of positive type on T, and that of the positive quadratic forms on \(E'\) .

Let K be a kernel of positive type on T, and q the associated quadratic form on \(E'\) . The Gaussian promeasure on E with variance q is also called the Gaussian promeasure on E with covariance K. If T is countable, Prop. 2 of No.~1 implies that this promeasure is a measure.

\begin{enumerate}
\def\labelenumi{\arabic{enumi})}
\setcounter{enumi}{2}
\tightlist
\item
  Let \(\mathbf{T}\) be a countable set. A kernel \(\delta\) on \(\mathbf{T}\) of positive type is defined by setting
\end{enumerate}

\[
\delta (t, t ^ {\prime}) = \left\{ \begin{array}{l l} 1 & \text { if } t = t ^ {\prime} \\ 0 & \text { if } t \neq t ^ {\prime}. \end{array} \right.\tag{24}
\]

The corresponding quadratic form is given by \(q\left(\sum_{t \in \mathbb{T}} c_t \varepsilon_t\right) = \sum_{t \in \mathbb{T}} c_t^2\) . For every \(t \in \mathbb{T}\) , let us denote by \(\mu_t\) the Gaussian measure on \(\mathbf{R}\) with variance 1;

one shows easily that the Gaussian measure on \(\mathbf{R}^{\mathrm{T}}\) with covariance \(\delta\) is equal to \(\bigotimes_{t\in T}\mu_t\) .

\begin{enumerate}
\def\labelenumi{\arabic{enumi})}
\setcounter{enumi}{3}
\tightlist
\item
  Let \(n \geqslant 1\) be an integer. A square matrix \(C = (c_{ij})\) of order \(n\) is said to be positive symmetric if it is symmetric and \(\sum_{i,j=1}^{n} c_{ij} x_i x_j \geqslant 0\) for any real \(x_1, \ldots, x_n\) ; it comes to the same to say that the mapping \((i,j) \mapsto c_{ij}\) is a kernel of positive type on the set \(\{1,2,\ldots,n\}\) . We may therefore speak of the Gaussian measure \(\gamma_C\) on \(\mathbf{R}^n\) , with covariance \(C\) ; it is characterized by the formula
\end{enumerate}

\[
\int_ {\mathbf {R} ^ {n}} e ^ {i (x _ {1} t _ {1} + \dots + x _ {n} t _ {n})} d \gamma_ {C} (t _ {1}, \dots , t _ {n}) = \exp \left(- \frac {1}{2} \sum_ {j, k = 1} ^ {n} c _ {j k} x _ {j} x _ {k}\right),\tag{25}
\]

for \(x_{1}, \ldots, x_{n}\) real. From Prop. 6 of No.~5 (formula (19)), one deduces

\[
\int_ {\mathbf {R} ^ {n}} t _ {j} t _ {k} d \gamma_ {C} (t _ {1}, \dots , t _ {n}) = c _ {j k} \quad (1 \leqslant j, k \leqslant n).\tag{26}
\]

From Prop. 5 of No.~5, one deduces the formula

\[
u (\gamma_ {C}) = \gamma_{U C\,{}^{t}U}  ,\tag{27}
\]

where u is a linear mapping of \(R^{n}\) into \(R^{m}\) with matrix U. Moreover, one sees easily (cf.~the first part of the proof of Prop. 4 of No.~5) that if \(I_{n}\) is the identity matrix of order n, then the measure \(\gamma_{I_{n}}\) admits the density

\[
(2 \pi) ^ {- n / 2} \exp \left(- \frac {1}{2} (t _ {1} ^ {2} + \dots + t _ {n} ^ {2})\right)
\]

with respect to the Lebesgue measure \(\lambda_{n}\) on \(\mathbf{R}^{n}\) .

We are going to show that if the matrix \(C\) is invertible, with inverse \(D = (d_{jk})\) , then

\[
\begin{array}{l} d \gamma_ {C} (t _ {1}, \dots , t _ {n}) = \\ (2 \pi) ^ {- n / 2} (\det D) ^ {1 / 2} \left(\exp \left(- \frac {1}{2} \sum_ {j, k = 1} ^ {n} d _ {j k} t _ {j} t _ {k}\right)\right) d t _ {1} \dots d t _ {n}. \end{array}\tag{28}
\]

For, if \(C\) is invertible, the quadratic form \(q\) on \(\mathbf{R}^n\) defined by

\[
q (x _ {1}, \dots , x _ {n}) = \sum_ {j, k = 1} ^ {n} c _ {j k} x _ {j} x _ {k}
\]

is nondegenerate. Using the existence of a basis of \(R^{n}\) orthonormal for q, one proves the existence of a square matrix U of order n such that \(C = U \cdot {}^{t}U\) , whence \(\gamma_{C} = u(\gamma_{I_{n}})\) by (27) (where u denotes the automorphism of \(R^{n}\) with matrix U). Let Q be the quadratic form on \(R^{n}\) defined by

\[
Q (t _ {1}, \dots , t _ {n}) = t _ {1} ^ {2} + \dots + t _ {n} ^ {2};
\]

then

\[
\gamma_ {I _ {n}} = (2 \pi) ^ {- n / 2} e ^ {- Q / 2} \cdot \lambda_ {n},
\]

whence

\[
u (\gamma_ {I _ {n}}) = (2 \pi) ^ {- n / 2} e ^ {- (\mathrm{Q} \circ u ^ {- 1}) / 2} \cdot u (\lambda_ {n}).\tag{29}
\]

It is immediate that the quadratic form \(\mathbf{Q} \circ u^{-1}\) on \(\mathbf{R}^n\) takes the value \(\sum_{j,k=1}^{n} d_{jk} t_j t_k\) at the point \((t_1, \ldots, t_n)\) , and Prop. 15 of Ch. VII, §1, No.~10 shows that

\[
u (\lambda_ {n}) = (\det U) ^ {- 1} \cdot \lambda_ {n} = (\det D) ^ {1 / 2} \cdot \lambda_ {n}.\tag{30}
\]

Formula (28) then follows from this.

